\lecture{15}{}{Operations with Matrices}
\section{Operations}
\begin{definition}[Matrix Addition]\label{def:matrix-addition}
    Let $A = [a_{ij}]_{m \times n}$ and $B = [b_{ij}]_{m \times n}$ be matrices with entries in the same field $F$. The \textbf{sum} $A + B$ is the matrix
    \[
    A + B \coloneqq [a_{ij} + b_{ij}]_{m \times n}
    \]
\end{definition}

\begin{note}
    Matrix addition is only defined for matrices of the same order. For example, the sum
    \[
    \begin{bmatrix} 1 & 2 \\ 3 & 4 \end{bmatrix} + \begin{bmatrix} 5 & 6 & 7 \\ 8 & 9 & 0 \end{bmatrix}
    \]
    is undefined.
\end{note}

\begin{proposition}[Properties of Matrix Addition]\label{prop:matrix-addition-properties}
    Matrix addition satisfies the following properties:
    \begin{enumerate}[label=(\roman*)]
        \item \textbf{Closure}: For all $A, B \in M_{m \times n}(F)$, $A + B \in M_{m \times n}(F)$
        \item \textbf{Commutativity}: For all $A, B \in M_{m \times n}(F)$, $A + B = B + A$
        \item \textbf{Associativity}: For all $A, B, C \in M_{m \times n}(F)$, $(A + B) + C = A + (B + C)$
        \item \textbf{Existence of neutral element}: There exists $O \in M_{m \times n}(F)$ such that for all $A \in M_{m \times n}(F)$, $A + O = A$
        \item \textbf{Existence of additive inverse}: For every $A \in M_{m \times n}(F)$, there exists $B \in M_{m \times n}(F)$ such that $A + B = O$
    \end{enumerate}
\end{proposition}

\begin{proof}
    Let $A = [a_{ij}]_{m \times n}$ and $B = [b_{ij}]_{m \times n}$ be matrices in $M_{m \times n}(F)$.
    
    \textbf{(i) Closure}: By definition, $A + B = [a_{ij} + b_{ij}]_{m \times n}$. Since $a_{ij}, b_{ij} \in F$ and $F$ is closed under addition, $a_{ij} + b_{ij} \in F$. Therefore $A + B \in M_{m \times n}(F)$.
    
    \textbf{(ii) Commutativity}: For all $i,j$, $[A + B]_{ij} = a_{ij} + b_{ij} = b_{ij} + a_{ij} = [B + A]_{ij}$ by commutativity of addition in $F$. Therefore $A + B = B + A$.
    
    \textbf{(iii) Associativity}: For all $i,j$, $[(A + B) + C]_{ij} = (a_{ij} + b_{ij}) + c_{ij} = a_{ij} + (b_{ij} + c_{ij}) = [A + (B + C)]_{ij}$ by associativity of addition in $F$. Therefore $(A + B) + C = A + (B + C)$.
    
    \textbf{(iv) Neutral element}: Let $O = [0]_{m \times n}$ be the matrix whose entries are all zeros of $F$. Then for all $i,j$, $[A + O]_{ij} = a_{ij} + 0 = a_{ij} = [A]_{ij}$. Therefore $A + O = A$.
    
    \textbf{(v) Additive inverse}: Let $B = [-a_{ij}]_{m \times n}$, where $-a_{ij}$ is the additive inverse of $a_{ij}$ in $F$. Then for all $i,j$, $[A + B]_{ij} = a_{ij} + (-a_{ij}) = 0 = [O]_{ij}$. Therefore $A + B = O$.
\end{proof}

\begin{note}
    The additive neutral element is called the \textbf{zero matrix}. As usual, associativity implies generalized associativity, so we can omit parentheses in sums.
\end{note}

\begin{definition}[Matrix Subtraction]\label{def:matrix-subtraction}
    Let $A, B \in M_{m \times n}(F)$. The \textbf{difference} is defined as
    \[
    A - B \coloneqq A + (-B)
    \]
    where $-B$ is the additive inverse of $B$.
\end{definition}

\begin{definition}[Scalar Multiplication]\label{def:scalar-multiplication}
    Let $A = [a_{ij}]_{m \times n} \in M_{m \times n}(F)$ and let $t \in F$. The \textbf{scalar multiple} $tA$ is the matrix
    \[
    tA \coloneqq [ta_{ij}]_{m \times n}
    \]
\end{definition}

\begin{theorem}[Properties of Scalar Multiplication]\label{thm:scalar-multiplication-properties}
    For any $A, B \in M_{m \times n}(F)$ and any $t, s \in F$:
    \begin{enumerate}[label=(\roman*)]
        \item $tA \in M_{m \times n}(F)$
        \item $s(tA) = (st)A$
        \item $(s + t)A = sA + tA$
        \item $t(A + B) = tA + tB$
        \item $1A = A$
        \item $0A = O_{m \times n}$
        \item $(-1)A = -A$
    \end{enumerate}
\end{theorem}

\begin{definition}[Matrix Multiplication]\label{def:matrix-multiplication}
    Let $A = [a_{ij}]_{m \times n}$ and $B = [b_{ij}]_{n \times q}$ be matrices with entries from the same field $F$. The \textbf{product} $AB$ is the $m \times q$ matrix defined by
    \[
    [AB]_{ij} = \sum_{k=1}^{n} a_{ik}b_{kj} = a_{i1}b_{1j} + a_{i2}b_{2j} + \cdots + a_{in}b_{nj}
    \]
\end{definition}

\begin{note}
    Key observations about matrix multiplication:
    \begin{itemize}
        \item The product $AB$ is defined when the number of columns in $A$ equals the number of rows in $B$
        \item The resulting matrix $AB$ has as many rows as $A$ and as many columns as $B$
        \item To compute $[AB]_{ij}$, multiply each entry in row $i$ of $A$ by the corresponding entry in column $j$ of $B$, then sum these $n$ products
        \item If $A$ is a $1 \times n$ row matrix and $B$ is an $n \times 1$ column matrix, then $AB$ is a $1 \times 1$ matrix (scalar)
        \item If $A$ is an $n \times 1$ column matrix and $B$ is a $1 \times n$ row matrix, then $BA$ is an $n \times n$ square matrix
        \item If $A$ is not square, then $AA$ is not defined
    \end{itemize}
\end{note}


\begin{note}
    Matrix multiplication is defined this way because it corresponds to composition of linear transformations. If $F$ is the matrix representing a linear function $f$ and $G$ represents $g$, then $FG$ represents the composition $f \circ g$, but only when this product is defined as in Definition~\ref{def:matrix-multiplication}.
\end{note}